\section{One-sided pointwise inversion and the physical denominator}
\label{sec:v92-lower}
Write $P=m^2p_{n,R}(k,m,t)$, $G=\mathcal L_{m,R}(k_1,k_2,t,n)$
and $f=P-G$. The positive remainder has a consequence in essential
supremum which does not require an upper bound for $P$. All suprema
in this section retain the original radius and exact integer labels.
The lower law itself is already Supplement
Theorem~\ref{supp-thm:v51-positive-error}, equation~\eqref{supp-eq:v51-lower-local-law}.
We restate its short proof to make the denominator and normalization
arguments self-contained; it is not a newly discovered lower law.

\begin{theorem}[The pointwise lower half of raw inversion]
\label{thm:v92-lower}
Under Input~\ref{input:v92-scalar}, for each fixed $M>0$,
\begin{equation}\label{eq:v92-lower}
 d_m(M):=\sup_{R,n,k}\esssup_{\mathcal C(M)}[G-P]_+
                                             \longrightarrow0.
\end{equation}
For any fixed $d>0$ put $D_m(d)=\mathcal C(M)\cap\{G\ge d\}$.
For all sufficiently large $m$, uniformly over the parameters,
\begin{equation}\label{eq:v92-denominator}
 P\ge G-d_m(M)\ge d/2\quad\hbox{a.e. on }D_m(d),\qquad
 \esssup_{D_m(d)}[G/P-1]_+\le 2d_m(M)/d.
\end{equation}
In particular the original conditional roof kernels have a positive
physical denominator on these sets. No ratio is assigned on $G=0$.
\end{theorem}
\begin{proof}
For each previously fixed reconstruction band $B$, the exact identity
$f=b_B+e_B$, with $b_B=m^2b^{\varepsilon(B)}\ge0$, gives
$[G-P]_+\le |e_B|$. Take the indicated supremum and collision limsup.
The result is at most $C_MB^{-1/192}$. Taking $B\to\infty$ proves
\eqref{eq:v92-lower}. This argument uses positivity, not weak-norm
convergence or a pointwise embedding. For large $m$, $d_m(M)<d/2$;
then $P\ge G-d_m(M)$ and
$[G/P-1]_+=[G-P]_+/P\le2d_m(M)/d$.
The coarea representatives and null sets are those of the original
source. The order remains $\sup_R\esssup_t$.
\end{proof}

The floor $d$ in this theorem is indispensable to its uniform ratio
statement. It is not replaced by the unquantified condition $G>0$.
The absolute lower bound \eqref{eq:v92-lower}, by contrast, retains
all arithmetic classes. It says nothing about the positive part of
$P-G$.

\begin{lemma}[A bounded quantity with a critical weak bound]
\label{lem:v92-bounded-critical}
Let $q>1$, $|J|=L<\infty$, $0\le u\le H$ and
$\mathcal N_q(u;J)\le A$. For $1\le p<q$,
\begin{equation}\label{eq:v92-bounded-subcritical}
 \|u\|_{L^p(J)}\le\left(\frac q{q-p}\right)^{1/p}
                                      A L^{1/p-1/q}.
\end{equation}
For $p>q$,
\begin{equation}\label{eq:v92-bounded-supercritical}
 \|u\|_{L^p(J)}\le\left(\frac p{p-q}\right)^{1/p}
                                      A^{q/p}H^{1-q/p}.
\end{equation}
When $0<A\le HL^{1/q}$,
\begin{equation}\label{eq:v92-bounded-endpoint}
 \|u\|_{L^q(J)}^q\le A^q
                     \left(1+q\log\frac{HL^{1/q}}A\right).
\end{equation}
Zero cases have their continuous interpretation; when $A>HL^{1/q}$,
one may replace $A$ by $HL^{1/q}$. In particular, bounded functions
converging to zero in $\mathcal N_q$ converge strongly in every finite
$L^p$.
\end{lemma}
\begin{proof}
For $\lambda>0$, testing on $\{u>\lambda\}$ gives
$|\{u>\lambda\}|\le\min\{L,(A/\lambda)^q\}$.
Integrate $p\lambda^{p-1}$ times this bound over $0<\lambda<H$,
splitting at $\lambda_0=AL^{-1/q}$. For $p<q$ extending the upper
integral to infinity gives \eqref{eq:v92-bounded-subcritical}.
For $p>q$, the exact expression is at most
$pA^qH^{p-q}/(p-q)$. At $p=q$ it is
$A^q+qA^q\log(H/\lambda_0)$. These are the claimed bounds.
The natural bound $\mathcal N_q(u;J)\le HL^{1/q}$ justifies the
replacement. All statements are for individual finite-count functions.
\end{proof}

\begin{theorem}[The reciprocal density, without an upper-height premise]
\label{thm:v92-reciprocal}
Assume also Input~\ref{input:v92-critical}. On $D_m(d)$ define
$r_m=|1-G/P|$, and set it to zero outside this set. For all sufficiently
large $m$,
\begin{equation}\label{eq:v92-reciprocal-bound}
 0\le r_m\le1,\qquad
 \mathcal N_Q(r_m;J)\le\frac2d
                    \mathcal N_Q(f\one_{\mathcal C(M)};J),
 \qquad Q=32/31.
\end{equation}
Consequently $r_m\to0$ in every finite $L^p$ on each fixed roof
window, uniformly in its location and the exact parameters. The same
conclusion holds on the complete central window in measure $dt/\sqrt m$.
No strong critical-norm conclusion for the \emph{unbounded raw error}
$f$ is used or inferred.
\end{theorem}
\begin{proof}
Equation~\eqref{eq:v92-denominator} gives $r_m\le2|f|/d$.
If $P\ge G$ then $r_m=1-G/P\le1$. If $P<G$ then
$r_m\le2d_m(M)/d\le1$ for large $m$. The weak bound follows by
positive domination. Apply Lemma~\ref{lem:v92-bounded-critical}
and the inherited unweakened critical convergence. For the diffusive
window, use the same pointwise inequalities after the change
$t=m\bar\tau_R+\sqrt m\,u$. The inherited critical norm there
already includes its Jacobian. The domain in $u$ is contained in
$[-M,M]$, so the bounded-function lemma applies on that fixed interval.
\end{proof}
